\section{Does Math RL Hurt General Ability?}
\label{sec:general}

A common concern with task-specific RL is that the policy gains on the
target distribution at the cost of forgetting general capabilities. We
check this on two general-purpose multiple-choice benchmarks:
\code{HellaSwag}~\cite{hellaswag} and \code{MMLU}~\cite{mmlu}.

\subsection{Numbers}
\Cref{tab:general} extracts the two general columns from
\cref{tab:main} and adds the standard error from
\emph{lm-eval-harness} for both benchmarks.

\begin{table}[h]
\centering
\small
\begin{tabular}{lcc}
\toprule
\textbf{Model stage} & \textbf{HellaSwag acc} & \textbf{MMLU acc (5-shot)} \\
\midrule
qwen base       & $0.5018$            & $0.6019 \pm 0.0039$ \\
qwen base sft   & $0.5068$            & $0.5993 \pm 0.0039$ \\
grpo sampled24k & $0.5101$            & $0.6030 \pm 0.0039$ \\
dapo            & $0.5105$            & $0.6014 \pm 0.0039$ \\
grpo            & $0.5112$            & $0.5998 \pm 0.0039$ \\
\bottomrule
\end{tabular}
\caption{General-ability evaluation. \code{HellaSwag} is 0-shot;
\code{MMLU} is 5-shot, aggregated across 57 subjects with 14\,042
questions, weighted by subject size. Standard error is reported by
\emph{lm-eval-harness}.}
\label{tab:general}
\end{table}

\subsection{Reading the table}
\paragraph{HellaSwag.}
The HellaSwag accuracy moves \emph{slightly upward} along the
training pipeline, from $0.502$ at the base to $0.511$ at the final
\code{grpo} checkpoint. The differences are small (less than $1.0$
absolute percentage point) but the direction is monotonic, suggesting
that the math pipeline at least does not damage commonsense
multiple-choice judgement.

\paragraph{MMLU.}
The MMLU column is the more demanding general test. All five model
stages land in the narrow band $[0.5993, 0.6030]$. With a per-row standard
error of $\pm 0.0039$, all five MMLU numbers are statistically
indistinguishable. The pipeline consequently leaves multi-subject knowledge
ability essentially unchanged --- neither the small SFT step nor the much
larger RL step measurably degrades it.

\subsection{Why this matters}
The combined picture from \cref{tab:main}, \cref{tab:general},
\cref{fig:resplen}, and \cref{fig:reward} is consistent with a clean
interpretation: a verifiable-reward RL stage on a 1.5B base \emph{can}
deliver multi-fold accuracy gains on math reasoning while keeping general
knowledge ability intact. A negative result on either general benchmark
would have suggested that the reward shape is too narrow or that the RL
update is over-fitting to surface-level math patterns. We do not see that
here.

\begin{remark}
The MMLU evaluation is intentionally run \emph{without} the math chat
template: the lm-eval-harness MMLU task scores by comparing the
loglikelihood of the four answer letters \code{A}/\code{B}/\code{C}/\code{D}
under the model, and a chat template would change the conditioning text
in a way that breaks comparability across the five models. The same
applies to \code{HellaSwag}.
\end{remark}
